\chapter{Full Mock Examination}
\label{app:mock-exam}

\section*{Instructions}

This mock examination contains 62 questions designed to simulate the actual IBM C1000-185 certification exam.

\textbf{Time Limit}: 90 minutes\\
\textbf{Passing Score}: 44/62 (71\%)\\
\textbf{Format}: Multiple choice, scenario-based questions

\textbf{Exam Tips}:
\begin{itemize}[leftmargin=*]
    \item Read each question carefully before selecting an answer
    \item Note keywords like "BEST", "MOST", "PRIMARY", "LEAST"
    \item Manage your time: aim for 90 seconds per question
    \item Mark difficult questions and return to them
    \item Don't change answers unless you're certain - first instinct is often correct
\end{itemize}

\textbf{Scoring}:
\begin{itemize}[leftmargin=*]
    \item 56-62 correct (90-100\%): Excellent - You're very well prepared
    \item 50-55 correct (80-89\%): Good - Review weak areas before exam
    \item 44-49 correct (71-79\%): Passing - More study recommended
    \item Below 44 (Below 71\%): More preparation needed - review thoroughly
\end{itemize}

\vspace{1cm}

\section{Section 1: Analyze and Design (Questions 1-10)}

\subsection*{Question 1}

A financial services company wants to build a chatbot that answers customer questions about investment products. The answers must be grounded in the company's regulatory-approved documentation. Which architecture pattern is MOST appropriate?

\textbf{A.} Simple prompt-based inference with a general foundation model\\
\textbf{B.} Fine-tuning a model on customer FAQs\\
\textbf{C.} Retrieval-Augmented Generation (RAG) with regulatory documents\\
\textbf{D.} Agent-based system with web search capabilities

\vspace{0.5cm}

\subsection*{Question 2}

Which of the following is a PRIMARY limitation of Large Language Models that would make them unsuitable for real-time stock trading recommendations?

\textbf{A.} Context window size\\
\textbf{B.} Knowledge cutoff date\\
\textbf{C.} Hallucination potential\\
\textbf{D.} Bias in training data

% Questions continue through Q62...

\section{Section 2: Prompt Engineering (Questions 11-20)}

\subsection*{Question 11}

When using few-shot prompting, what is the RECOMMENDED number of examples to include for optimal results?

\textbf{A.} 1-2 examples\\
\textbf{B.} 3-5 examples\\
\textbf{C.} 10-15 examples\\
\textbf{D.} 20+ examples

% More questions...

\section{Section 3: Fine-Tuning (Questions 21-40)}

% Questions on LoRA, InstructLab, dataset preparation, etc.

\section{Section 4: RAG (Questions 41-51)}

% Questions on embeddings, vector databases, LangChain

\section{Section 5: Deployment (Questions 52-59)}

% Questions on deployment strategies, versioning, monitoring

\section{Section 6: Integration (Questions 60-62)}

% Questions on API integration, orchestration

\vspace{2cm}

\textbf{END OF EXAMINATION}

\textit{Answers and detailed explanations are provided in Appendix C.}

\cleardoublepage
